Spatial count data with excess zeros arise frequently in regional science, epidemiology, and urban economics, yet the existing econometric toolkit lacks a unified frequentist framework that simultaneously addresses spatial spillovers through a structural autoregressive mechanism, zero-inflation, overdispersion, and unobserved heterogeneity.
We introduce the Zero-Inflated Spatial Autoregressive (ZISAR) model, which places a simultaneous autoregressive (SAR) structure on the latent linear predictor of a zero-inflated negative binomial specification: the log-conditional mean follows $\boldsymbol{\eta} = (\mathbf{I} - \rho\mathbf{W})^{-1}(\mathbf{X}\boldsymbol{\beta} + \mathbf{u})$, where the spatial multiplier $(\mathbf{I} - \rho\mathbf{W})^{-1}$ generates endogenous feedback effects among neighbouring units and unit-level Gaussian random effects $\mathbf{u}$ capture residual heterogeneity.
The general formulation allows a separate SAR process in the zero-inflation component, with a restricted version confining spatial dependence to the count process.
Estimation proceeds by maximising the Laplace-approximated marginal likelihood, implemented via Template Model Builder (\texttt{TMB}), which exploits automatic differentiation and sparse-matrix algebra for scalable inference.
We establish consistency and asymptotic normality of the estimator and derive closed-form expressions for direct, indirect, and total marginal effects that account for both the mixture structure and the spatial multiplier.
For large $n$, the log-determinant of the spatial filter is obtained from a sparse Cholesky factorisation, so that no dense $n\times n$ matrix is formed.
A Monte~Carlo study demonstrates that ignoring either zero-inflation or spatial dependence induces substantial bias in the autoregressive and regression coefficients, whereas ZISAR recovers the data-generating parameters with low bias and an RMSE that decreases with the sample size, and that residual-based Moran tests can miss spatial dependence that the likelihood detects.
An application to domestic solar photovoltaic installations in the 32{,}322 small areas of England after the 2016 reform of the Feed-in Tariff, where a third of the areas record no installation, estimates strong spatial dependence in the latent propensity to adopt ($\hat\rho = 0.63$), attributes about one zero in ten to the structural regime of flat-dominated urban cores, and shows that aspatial models misread the spatial clustering of adoption as effects of tenure and age.
Companion \texttt{R} code and the \texttt{TMB} template implementing the estimator are provided.
